\chapter{Conclusiones y trabajo futuro}
\label{cap:conclusiones}

\section{Conclusiones}

El trabajo permitió construir y verificar una herramienta que parte de una fotografía real, registra la hoja impresa, segmenta los toppers con SAM 2 y transforma las siluetas en geometría DXF. La propuesta responde al problema observado en Innokey desde un cambio sencillo de proceso: laminar primero la hoja completa sobre el soporte y calcular el corte después, cuando la impresión y el material ya ocupan su posición definitiva.

La primera conclusión es que la inteligencia artificial resulta viable como parte operativa del sistema, aunque no haya sido el método con mayor puntuación real. SAM 2 alcanzó un IoU de $0.9536$ sobre las trece capturas y conservó los ocho objetos en todas ellas. Su máscara alimentó la extracción de contornos y la exportación. La prueba sobre la captura 20 terminó con ocho contornos y un DXF válido. Por tanto, la integración de IA no es una comparación aislada ni una opción decorativa de la interfaz.

La segunda conclusión es que un método clásico bien adaptado puede seguir siendo la referencia más fuerte en un montaje controlado. Su IoU real de $0.9787$ superó al de SAM 2. El fondo claro, los marcadores conocidos y la cantidad fija de figuras favorecen reglas de color, borde y morfología. Este resultado no se ocultó ni se forzó una ventaja artificial del modelo neuronal. Documentar que la IA no siempre supera a una solución específica forma parte del aporte académico.

La tercera conclusión surge de Random Forest. El modelo seleccionado logró un IoU sintético de $0.9902$, el mayor de la prueba, pero descendió a $0.2877$ en real. La optimización también mostró que el Jaccard calculado sobre píxeles muestreados puede favorecer un candidato que luego fragmenta la hoja completa. La validación por lámina evitó aceptar ese resultado. En este caso, la brecha sintético-real no es una explicación genérica, sino una diferencia cuantificada de $0.7025$ puntos de IoU.

La cuarta conclusión se refiere al registro. Las desviaciones estándar medias de $0.289\,\text{mm}$ en X y $0.251\,\text{mm}$ en Y indican que la homografía y el WCS proporcionan una base computacional repetible. La deriva máxima promedio de $0.820\,\text{mm}$ también advierte que el error angular crece con la distancia al origen. El resultado es suficiente para justificar la continuación del prototipo, pero no sustituye una medición física del corte.

\section{Cumplimiento de los objetivos}

La Tabla~\ref{tab:cumplimiento_objetivos} relaciona cada objetivo específico con la evidencia obtenida. El objetivo general se considera cumplido en el alcance software y computacional definido. La reducción efectiva de reprocesos en producción permanece como una hipótesis de impacto que deberá medirse mediante una campaña física.

\begin{table}[htbp]
\centering
\small
\caption{Cumplimiento de los objetivos específicos.}
\label{tab:cumplimiento_objetivos}
\begin{tabularx}{\textwidth}{p{1.6cm}Yp{1.8cm}}
\toprule
\textbf{Objetivo} & \textbf{Evidencia principal} & \textbf{Estado} \\
\midrule
OE1 & Procedimiento de captura con teléfono, soporte removible, hoja A4 y trece fotografías bajo variaciones controladas. & Cumplido \\
OE2 & Detección de cuatro marcadores, rectificación a $2100\times2970$ píxeles y WCS válido en diez capturas. & Cumplido \\
OE3 & Banco de 48 hojas, manifiesto canónico y partición de 36 hojas de entrenamiento, 6 de validación y 6 de prueba. & Cumplido \\
OE4 & Búsqueda aleatoria, 40 ajustes con validación agrupada y selección final de RF v1 con umbral 0.70 sobre hojas completas. & Cumplido \\
OE5 & Comparación de cinco prompts, selección de caja con margen del 5\%, evaluación sintética y real e integración de SAM 2 en la GUI. & Cumplido \\
OE6 & Extracción y simplificación de contornos, transformación a milímetros y exportación DXF protegida por WCS. & Cumplido \\
OE7 & Comparación mediante IoU, Dice, Boundary F1, componentes y tiempo, con referencia real asistida y limitaciones declaradas. & Cumplido \\
\bottomrule
\end{tabularx}
\end{table}

\section{Aportes del trabajo}

El TFM deja cuatro aportes concretos. El primero es una arquitectura híbrida donde la geometría clásica registra la escena y un modelo fundacional produce la máscara de uso operativo. El segundo es un protocolo que separa entrenamiento, validación y prueba por hoja y que selecciona Random Forest con una medida coherente con la vectorización. El tercero es una comparación sobre dos dominios con una referencia real que ya no toma la salida clásica como verdad de píxel. El cuarto es una implementación reproducible, acompañada de manifiestos, checksums, pruebas y resultados tabulares.

El hallazgo principal puede resumirse sin contradicciones. Random Forest fue sobresaliente dentro del dominio que aprendió y falló al transferirse. SAM 2 fue menos preciso que Random Forest en sintético, pero mucho más estable en real. La visión clásica obtuvo el mejor resultado real. Aun así, SAM 2 demostró calidad suficiente para generar la salida de la aplicación. Cada método aportó una respuesta distinta al problema y la decisión final se apoyó en datos.

\section{Trabajo futuro}

La prioridad siguiente es ampliar el conjunto real. Se requieren varias hojas físicas, diseños no observados, materiales diferentes y anotaciones manuales independientes. Ese banco permitiría estimar generalización y ajustar SAM 2 sin convertir las vistas repetidas de una sola hoja en ejemplos aparentemente independientes.

La segunda línea es mejorar la geometría. Una calibración intrínseca de la cámara reduciría la distorsión antes de la homografía. También conviene aumentar la longitud de la marca WCS y estudiar referencias distribuidas para disminuir la sensibilidad angular en la zona más alejada.

La tercera línea se centra en la IA. Un detector ligero podría sustituir el localizador clásico de cajas y permitir una evaluación neuronal más autónoma. Con un conjunto real suficiente, sería razonable comparar ajuste fino de SAM 2, adaptadores de bajo rango y generación de datos sintéticos más cercana a la cámara. Random Forest podría revisarse con una función de optimización calculada directamente sobre hojas completas, aunque su coste limitaría la cantidad de candidatos.

La cuarta línea es productiva. Debe ejecutarse una campaña de cortes sobre MDF y acrílico, medir el error entre borde impreso y borde cortado, estimar el ancho real del haz y ajustar la compensación. Solo entonces podrá cuantificarse la reducción de reprocesos, el tiempo de preparación y la aceptación visual por parte del taller.

Finalmente, la aplicación puede optimizarse para uso cotidiano. La inferencia en GPU, la caché del embedding de imagen y un instalador con verificación automática del checkpoint reducirían la fricción. Estas mejoras no cambian el resultado central del TFM, pero acercan el prototipo validado a una herramienta de producción.
